% Part: first-order-logic
% Chapter: completeness
% Section: outline

\documentclass[../../../include/open-logic-section]{subfiles}

\begin{document}

\iftag{FOL}
      {\olfileid{fol}{com}{out}}
      {\olfileid{pl}{com}{out}}

\olsection{સાબિતીની રૂપરેખા}

પૂર્ણતા પ્રમેયની સાબિતી થોડી જટિલ છે અને પહેલી વાર વાંચતી વખતે તેમાં
અટવાઈ જવું સહેલું છે. તેથી તેની રૂપરેખા જોઈ લઈએ. પહેલું પગલું
દૃષ્ટિકોણ બદલવાનું છે; આ બદલાવ આપણને સાબિતી તરફનો માર્ગ દેખાડે છે.
પૂર્ણતાને ``જ્યારે પણ $\Gamma \Entails !A$ હોય ત્યારે $\Gamma \Proves
!A$ હોય'' એમ વિચારીએ, ત્યારે વિચાર સૂઝવો પણ અઘરો થઈ શકે છે: $\Gamma
\Proves !A$ બતાવવા માટે આપણે !!a{derivation} શોધવાની છે, અને $\Gamma
\Entails !A$ એ પૂર્વધારણા તેમાં કોઈ રીતે મદદ કરતી હોય તેમ દેખાતું નથી.
કેટલીક સાબિતી-પદ્ધતિઓ માટે !!a{derivation} સીધી રચી શકાય છે, પરંતુ
આપણે સહેજ જુદો અભિગમ લઈશું. જરૂરી દૃષ્ટિકોણ-પરિવર્તન આ છે: પૂર્ણતાનું
નિરૂપણ ``જો $\Gamma$ સુસંગત હોય, તો તે સંતોષ્ય છે'' એમ પણ કરી શકાય.
કદાચ~$\Gamma$માંની માહિતી અને તે સુસંગત છે એ પૂર્વધારણા સાથે વાપરીને
આપણે~$\Gamma$ના દરેક \iftag{FOL}{!!{sentence}}{!!{formula}}ને સંતોષતી
\iftag{FOL}{!!a{structure}}{!!a{valuation}} રચી શકીએ. આખરે, આપણે કેવી
\iftag{FOL}{!!{structure}}{!!{valuation}} શોધીએ છીએ તે જાણીએ છીએ:
~$\Gamma$ જેવું તેનું વર્ણન કરે છે એવી!{}

જો~$\Gamma$માં માત્ર \iftag{FOL}{પરમાણુક
  !!{sentence}s}{!!{propositional variable}s} હોય, તો તેનો નિદર્શ રચવો
સહેલો છે.\iftag{FOL}{ ધારો કે બધાં પરમાણુક !!{sentence}s
  $\Atom{P}{a_1,\dots,a_n}$ સ્વરૂપનાં છે, જ્યાં $a_i$ !!{constant}s
  છે.}{} આપણે માત્ર એટલું કરવાનું છે કે
\iftag{FOL}{%
  !!a{domain}~$\Domain{M}$ અને~$P$ માટે એવી નિયુક્તિ આપવી કે
  $\Sat{M}{\Atom{P}{a_1,\ldots,a_n}}$. પણ એ બહુ અઘરું નથી:
  $\Domain{M} = \Nat$, $\Assign{\Obj c_i}{M} = i$ મૂકો; અને દરેક
  $\Atom{P}{a_1,\ldots,a_n} \in \Gamma$ માટે ટપલ
  $\tuple{k_1,\dots,k_n}$ને~$\Assign{P}{M}$માં મૂકો, જ્યાં $k_i$ એ
  અચળ સંકેત~$a_i$નો સૂચક છે (એટલે કે $a_i \ident \Obj c_{k_i}$).}{%
  એવી !!a{valuation}~$\pAssign{v}$ કે દરેક $p \in \Gamma$ માટે
  $\pSat{v}{p}$. એટલે $\pAssign{v}(p)=\True$ જો અને માત્ર જો
  $p\in\Gamma$ એમ મૂકો.}

હવે ધારો કે~$\Gamma$માં કોઈ !!{formula}~$\lnot !B$ છે, જેમાં $!B$
પરમાણુક છે. $\iftag{FOL}{\Struct{M}}{\pAssign{v}}$ની રચના $\lnot !B$ને
સાચું બનાવવાની શક્યતામાં દખલ કરે છે એવી ચિંતા થઈ શકે. પણ અહીં જ
~$\Gamma$ની સુસંગતતા કામ આવે છે: જો $\lnot !B \in \Gamma$ હોય, તો
$!B \notin \Gamma$; નહિ તો~$\Gamma$ અસુસંગત હોત. અને જો $!B \notin
\Gamma$ હોય, તો આપણી~$\iftag{FOL}{\Struct{M}}{\pAssign{v}}$ની રચના મુજબ
$\iftag{FOL}{\Sat/{M}}{\pSat/{v}}{!B}$, તેથી
$\iftag{FOL}{\Sat{M}}{\pSat{v}}{\lnot !B}$. અત્યાર સુધી બધું બરાબર છે.

જો~$\Gamma$માં જટિલ, બિન-પરમાણુક સૂત્રો હોય તો શું? માનો કે તેમાં
$!A \land !B$ છે. તેને સાચું બનાવવા માટે આપણે~$\Gamma$માં $!A$ અને
$!B$ બંને હોય એ રીતે આગળ વધવું જોઈએ. અને જો $!A \lor !B \in \Gamma$
હોય, તો તેમાંથી ઓછામાં ઓછું એક સાચું બનાવવું પડશે, એટલે કે તેમાંથી એક
~$\Gamma$માં હોય એ રીતે આગળ વધવું પડશે.

આમાંથી નીચેનો વિચાર સૂઝે છે: આપણે~$\Gamma$માં વધુ !!{formula}s ઉમેરીએ,
જેથી (a)~મળતો ગણ સુસંગત રહે; (b)~દરેક શક્ય !!{sentence}~$!A$ માટે
મળતા ગણમાં કાં તો~$!A$ હોય અથવા $\lnot !A$ હોય; અને (c)~જ્યારે પણ
ગણમાં $!A \land !B$ હોય ત્યારે તેમાં $!A$ અને $!B$ બંને હોય, જ્યારે
ગણમાં $!A \lor !B$ હોય ત્યારે તેમાં $!A$ અથવા $!B$માંથી ઓછામાં ઓછું એક
પણ હોય, વગેરે. આપણે આ પ્રક્રિયા ચાલુ રાખીએ (જરૂર પડે તો અનંત સુધી).
આ રીતે ઉમેરાયેલાં બધાં !!{formula}sના ગણને~$\Gamma^*$ કહો. પછી ઉપરની
રચના એવી \iftag{FOL}{!!a{structure}~$\Struct{M}$}{!!a{valuation}~$\pAssign{v}$}
આપશે કે તે~$\Gamma^*$નાં બધાં વાક્યોને સંતોષે છે એવું આપણે અનુમાનથી
સાબિત કરી શકીશું; અને $\Gamma \subseteq \Gamma^*$ હોવાથી તે
~$\Gamma$નાં બધાં વાક્યોને પણ સંતોષશે. ખરેખર, (a) અને~(b)ની ખાતરી જ
પૂરતી હોવાનું બહાર આવે છે. જે વાક્યોના ગણ માટે~(b) સાચું હોય તેને
\emph{પૂર્ણ} કહે છે. તેથી સુસંગત ગણ~$\Gamma$ને સુસંગત અને પૂર્ણ ગણ
~$\Gamma^*$ સુધી વિસ્તારવાનું આપણું કાર્ય હશે.

\sourcecorrection{OLCO-003}{મૂળની \texttt{outline.tex}, પંક્તિ~68માં
શરત~(b)ને માત્ર ``પરમાણુક'' વાક્યો સુધી મર્યાદિત કરી છે, પરંતુ
પંક્તિઓ~78--80ની પૂર્ણ ગણની વ્યાખ્યા દરેક વાક્ય માટે નિર્ણય માગે છે.
અહીં શરત~(b)માં ``પરમાણુક'' મર્યાદા દૂર કરી છે.}

\sourcecorrection{OLCO-004}{મૂળની \texttt{outline.tex}, પંક્તિ~76માં
``all sentence'' લખાયું છે. અહીં બહુવચન અર્થ અનુસાર ``બધાં વાક્યો''
લખ્યું છે.}

\iftag{FOL}{%
આ યોજનામાં એક અડચણ છે: જો $\lexists[x][!A(x)] \in \Gamma$ હોય, તો
આપણે કોઈ !!{constant}~$c$ પસંદ કરીને આ પ્રક્રિયામાં $!A(c)$ ઉમેરવા
ઇચ્છીએ. પરંતુ એવું હંમેશાં કરી શકાય છે તે કેવી રીતે જાણીએ? શક્ય છે કે
આપણી ભાષામાં બહુ થોડાં !!{constant}s હોય અને તેમાંના દરેક માટે
$\lnot !A(c) \in \Gamma$ હોય. આપણે સાથે $!A(c)$ પણ ઉમેરી શકીએ નહિ,
કારણ કે તેનાથી ગણ અસુસંગત થઈ જાય; અને પછી~$\Struct{M}$એ $!A(c)$ સાચું
બનાવવાનું છે કે $\lnot !A(c)$ તે જાણી શકાય નહિ. વધુમાં, એવું પણ બની
શકે કે~$\Gamma$માં એવી ભાષાનાં જ વાક્યો હોય જેમાં એકેય અચળ સંકેત ન હોય
(દાખલા તરીકે, ગણસિદ્ધાંતની ભાષા).

આ સમસ્યાનો ઉકેલ એટલો જ છે કે શરૂઆતમાં અનંત સંખ્યામાં અચળો અને તેમને
પરિમાણકો સાથે યોગ્ય રીતે જોડતાં વાક્યો ઉમેરી દઈએ. (અલબત્ત, તેનાથી
અસુસંગતતા પેદા ન થઈ શકે તેની ચકાસણી કરવી પડશે.)

પરમાણુક વાક્યોમાં માત્ર !!{constant}s હોય ત્યારે આપણી મૂળ રચના સારી
રીતે કામ કરે છે. પરંતુ ભાષામાં !!{function}s પણ હોઈ શકે. એ કિસ્સામાં
બધું યોગ્ય બને તે માટે આ !!{function}sને નિયુક્ત કરવા યોગ્ય
~$\Nat$-પરનાં વિધેયો શોધવા અઘરા હોઈ શકે. તેથી બીજી યુક્તિ વાપરીએ:
$\Obj c_i$નું અર્થઘટન કરવા~$i$ વાપરવાને બદલે, !!{constant}sના ગણને જ
વ્યાપકક્ષેત્ર લઈએ. પછી~$\Struct M$ દરેક !!{constant}ને પોતાને જ નિયુક્ત
કરી શકે: $\Assign{\Obj c_i}{M}=\Obj c_i$. પરંતુ અહીં અટકવું શા માટે?
$\Domain{M}$ને ભાષાનાં બધાં \emph{પદો}નો ગણ લઈએ!{} એમ કરીએ તો
!!{function}sને વિધેયો (જે પદોને દલીલો તરીકે લે અને પદોને જ મૂલ્યો
તરીકે આપે) નિયુક્ત કરવાની સ્પષ્ટ રીત મળે છે: $n$ પદો $t_1$, \dots,
~$t_n$ નિવેશ તરીકે મળે ત્યારે પદ $\Obj f^n_i(t_1,\dots,t_n)$ને મૂલ્ય
તરીકે આપતું વિધેય !!{function}~$\Obj f^n_i$ને નિયુક્ત કરીએ.

કોયડાનો છેલ્લો ભાગ~$\eq$ને કેવી રીતે સંભાળવો તે છે. !!{predicate}
~$\eq$નું અર્થઘટન નિશ્ચિત છે: $\Sat{M}{\eq[t][t']}$ જો અને માત્ર જો
$\Value{t}{M}=\Value{t'}{M}$. હવે પદ~$t$નું !!{value} એ~$t$ પોતે જ બને
એવી ગોઠવણી કરીએ, તો $t$ અને $t'$ એક જ પદ હોય તે સિવાય આ !!{structure}
$\eq[t][t']$ સ્વરૂપનું \emph{એકેય} વાક્ય સાચું બનાવશે નહિ. અને આ
નિશ્ચિતપણે સમસ્યા છે, કારણ કે !!{function}sવાળી ભાષામાંના લગભગ દરેક
રસપ્રદ સિદ્ધાંતમાં $t$ અને $t'$ એક જ પદ ન હોય એવાં વાક્યો
$\eq[t][t']$ પ્રમેય તરીકે હશે (દાખલા તરીકે, અંકગણિતના સિદ્ધાંતોમાં
$\eq[(\Obj 0+\Obj 0)][\Obj 0]$). આ સમસ્યા ઉકેલવા આપણે~$\Struct M$નું
વ્યાપકક્ષેત્ર બદલીએ: $\Domain{M}$માં પદોનો પદાર્થો તરીકે ઉપયોગ કરવાને
બદલે, પદોના ગણો વાપરીએ; દરેક ગણ એવો હોય કે~$\Gamma^*$નાં વાક્યો જે
પદોને સમાન હોવા માગે છે તે બધાં પદો તેમાં આવે. તેથી, દાખલા તરીકે,
જો~$\Gamma$ અંકગણિતનો સિદ્ધાંત હોય, તો આ ગણોમાંના એકમાં $\Obj 0$,
$(\Obj 0+\Obj 0)$, $(\Obj 0\times\Obj 0)$ વગેરે આવશે. આ ગણને આપણે
$\Obj 0$ને નિયુક્ત કરીશું, અને બહાર આવશે કે આ ગણ તેમાંનાં બધાં પદોનું
મૂલ્ય પણ છે; દાખલા તરીકે, $(\Obj 0+\Obj 0)$નું પણ. તેથી સુધારેલી
!!{structure}માં વાક્ય $\eq[(\Obj 0+\Obj 0)][\Obj 0]$ સાચું હશે.}{}

\sourcecorrection{OLCO-005}{મૂળની \texttt{outline.tex}, પંક્તિ~126માં
ભાગફલ રચનાના ગણો મૂળ~$\Gamma$નાં વાક્યો માગે તે પદોને જોડે છે એમ
કહ્યું છે. પંક્તિઓ~157--168 અને પછીની ઔપચારિક રચનામાં આ સંબંધ
પૂર્ણ વિસ્તાર~$\Gamma^*$ પરથી વ્યાખ્યાયિત થાય છે. અહીં તેથી
$\Gamma^*$ લખ્યું છે.}

હવે આપણે શું કરીશું તે જોઈએ. પહેલાં !!{complete} સુસંગત ગણોના ગુણધર્મો
તપાસીશું. ખાસ કરીને, !!a{complete} સુસંગત ગણમાં $!A \land !B$ હોય જો
અને માત્ર જો તેમાં $!A$ અને~$!B$ બંને હોય; તેમાં $!A \lor !B$ હોય જો
અને માત્ર જો તેમાંથી ઓછામાં ઓછું એક હોય; વગેરે સાબિત કરીશું
(\olref[ccs]{prop:ccs}).\iftag{FOL}{ પછી વાક્યોના ``સંતૃપ્ત'' ગણો
  વ્યાખ્યાયિત કરીને તપાસીશું. સંતૃપ્ત ગણમાં દરેક પરિમાણિત
  !!{sentence}ને તેના દૃષ્ટાંતો સાથે જોડતી શરતી વિધિઓ હોય છે
  (\olref[hen]{defn:henkin-exp}). દરેક સુસંગત ગણ~$\Gamma$ને હંમેશાં
  સંતૃપ્ત ગણ~$\Gamma'$ સુધી વિસ્તારી શકાય છે તેમ બતાવીશું
  (\olref[hen]{lem:henkin}). કોઈ ગણ સુસંગત, સંતૃપ્ત અને !!{complete}
  હોય, તો તેમાં એ ગુણધર્મ પણ હોય છે કે
  \iftag{prvEx}{$\lexists[x][!A(x)]$ તેમાં હોય જો અને માત્ર જો કોઈ બંધ
    પદ~$t$ માટે $!A(t)$ તેમાં હોય}{}\iftag{defEx,defAll}{}{ અને
  }\iftag{prvAll}{$\lforall[x][!A(x)]$ તેમાં હોય જો અને માત્ર જો દરેક
    બંધ પદ~$t$ માટે $!A(t)$ તેમાં હોય}{}
  (\olref[hen]{prop:saturated-instances}).}{}
પછી આપણે \iftag{FOL}{સંતૃપ્ત}{} સુસંગત ગણ
~$\iftag{FOL}{\Gamma'}{\Gamma}$ને લઈશું અને તેને
\iftag{FOL}{સંતૃપ્ત, સુસંગત અને}{સુસંગત અને} !!{complete} ગણ
~$\Gamma^*$ સુધી વિસ્તારી શકાય છે તેમ બતાવીશું
(\olref[lin]{lem:lindenbaum}). આ~$\Gamma^*$ વડે આપણે આપણો
\iftag{FOL}{પદ-નિદર્શ~$\Struct M(\Gamma^*)$}{!!{valuation}~$\pAssign v(\Gamma^*)$}
વ્યાખ્યાયિત કરીશું. ~ $\Gamma^*$માંનાં
\iftag{FOL}{પરમાણુક !!{sentence}s પદ-નિદર્શના !!{predicate}sનું
  અર્થઘટન નક્કી કરે છે}{!!{propositional
    variable}s સત્યમૂલ્ય-નિયુક્તિ નક્કી કરે છે}
(\olref[mod]{defn:termmodel}). પછી
\iftag{FOL}{સંતૃપ્ત,}{} પૂર્ણ સુસંગત ગણોના ગુણધર્મો વાપરીને ખરેખર
$\iftag{FOL}{\Sat{M(\Gamma^*)}}{\pSat{v(\Gamma^*)}}{!A}$ જો અને માત્ર જો
$!A\in\Gamma^*$ એવું બતાવીશું (\olref[mod]{lem:truth}), અને તેથી ખાસ
કરીને $\iftag{FOL}{\Sat{M(\Gamma^*)}}{\pSat{v(\Gamma^*)}}{\Gamma}$.
\iftag{FOL}{અંતે, જો~$\Gamma$માં~$\eq$ પણ હોય તો પદ-નિદર્શ કેવી રીતે
  વ્યાખ્યાયિત કરવો તે જોઈશું (\olref[ide]{defn:term-model-factor}) અને
  તે~$\Gamma^*$ને સંતોષે છે તેમ બતાવીશું (\olref[ide]{lem:truth}).}{}

\end{document}
